% !TEX root = userguide.tex

\chapter{Subdivision of simplexes (triangles, tetrahedrons)}
\pgst{tfemchA}
\label{chap:subdivide}

\def\chaptertitle{Subdivision of simplexes (triangles, tetrahedrons)}
\def\headdate{10th October 2012}

(This appendix has been contributed by Nick Jaensson of the TU/e).\bigskip

The subdivision of a square (the reference region in the eltree for quads) and a cube (the reference region in the 
eltree for hexahedrons) is quite straightforward. A square can be divided into four subsquares while a cube can be 
divided into eight subcubes. In both cases the edges of the subfigures will be half of the length of the edges of the 
initial simplex. However, for triangles and especially tetrahedrons, this subdivision is less trivial. 

\section{Subdivision in 2D: triangles}
In the eltree for quads, every subelement can be found by two coordinates (the bottom left and top right, given by 
\texttt{coor(1)} and \texttt{coor(2)} resp.). This system can be used for triangles by adding one piece of 
information: each triangle can be the upper of lower part of the square. The variable \texttt{conf} (for configuration)
contains this information by taking the value 1 for lower triangles and the value 2 for upper triangles as shown in 
in Figure \ref{fig:triangle_confs}.

\begin{figure}[bhtp!]
\begin{center}
\includegraphics[width = 0.3\linewidth]{confs}
\caption{There are two possible configurations for a triangle. }
\label{fig:triangle_confs}
\end{center}
\end{figure}

The subdivision is done by connecting the midpoints of the three edges, this will yield 3 triangles that are similar 
to the original one, and one that triangle that is flipped. The initial (or 'root' level) of the subdivision will 
take place on a triangle with \texttt{conf} = 1, \texttt{coor(1) = (0,0)} and \texttt{coor(2) = (1,1)}, which is 
equal to the reference region for the isoparametric elements in \tfem. The subdivision of both configurations with 
the resulting \texttt{conf} and \texttt{coor(1)} are shown in Figures \ref{fig:el1} and \ref{fig:el2}. Here, the 
italic numbers denote the configuration number, the encircled numbers give the numbering of the subelements, and the 
numbers in a square show the numbering of the vertices. Note that \texttt{coor(2)} is always equal to 
\texttt{coor(1) + h/2}.

\renewcommand{\arraystretch}{1.5}
\begin{figure}[bhtp!]
  \begin{minipage}[c]{0.45\linewidth}\centering
    \includegraphics[width=0.8\textwidth]{el1_triangle}
  \end{minipage}
  \begin{minipage}[c]{0.45\linewidth}\centering
    \begin{tabular}{c | p{4cm} | c }
%      \hline
      subel. & coor(1)  & conf  \\ 
      \hline%\hline
      1 & ($x_0$,     $y_0$)                     & 1 \\ %\hline
      2 & ($x_0$,    $y_0$)                     & 2 \\ %\hline
      3 & ($x_0$ + $\frac{h(1)}{2}$,  $y_0$)    & 1 \\ %\hline
      4 & ($x_0$, $y_0$ + $\frac{h(2)}{2}$ )    & 1 %\\
     % \hline
    \end{tabular}
  \end{minipage}
\caption{Subdivision of triangles with \texttt{conf = 1}.}    
\label{fig:el1}
\end{figure}


\renewcommand{\arraystretch}{1.5}
\begin{figure}[bhtp!]
  \begin{minipage}[c]{0.45\linewidth}\centering
    \includegraphics[width=0.8\textwidth]{el2_triangle}
  \end{minipage}
  \begin{minipage}[c]{0.45\linewidth}\centering
    \begin{tabular}{c | p{4cm} | c }
      %\hline
      subel. & coor(1)  & conf  \\ 
      \hline%\hline
      1 & ($x_0$ + $\frac{h(1)}{2}$,  $y_0$)                     & 2 \\% \hline
      2 & ($x_0$,    $y_0$ + $\frac{h(2)}{2}$)                     & 2 \\% \hline
      3 & ($x_0$ + $\frac{h(1)}{2}$,  $y_0$ + $\frac{h(2)}{2}$)    & 1 \\% \hline
      4 & ($x_0$ + $\frac{h(1)}{2}$, $y_0$ + $\frac{h(2)}{2}$ )    & 2 %\\
     % \hline
    \end{tabular}
  \end{minipage}
\caption{Subdivision of triangles with \texttt{conf = 2}.}    
\label{fig:el2}
\end{figure}

To fill the root element for triangles correctly, the first call to \texttt{fill\_node\_eltree} will always look 
like: \newline

\noindent
\texttt{coor(1,:) = [ 0.0\_dp, 0.0\_dp ] } \newline
\texttt{coor(2,:) = [ 1.0\_dp, 1.0\_dp ] } \newline
\texttt{conf = 1}\newline
\texttt{eltype = 1}\newline
\texttt{call fill\_node\_eltree ( eltree, coor, conf, eltype )}\newline





\section{Subdivision in 3D: tetrahedrons}
For tetrahedrons, a subdivision into congruent subtetrahedrons is not obvious. In fact, \cite{matousek2011} mention 
that the 'Hill-simplex' is the only tetrahedron that can be subdivided into eight congruent subtetrahedrons. A Hill-simplex
is constructed by the vertices (0,0,0), (1,0,0), (1,1,0) and (1,1,1). By performing the subdivision into eight subelements 
(as shown in \cite{matousek2011}), 5 configurations are found. These configurations can also be subdivided, after which
it becomes clear that there are only 6 returning configurations. The subdivisions for the 6 possible configurations are 
shown in Figures \ref{fig:tet_el1} to \ref{fig:tet_el6}. Here the italic numbers denote the configuration number, the 
encircled numbers give the numbering of the subelements, and the numbers in a square show the numbering of the 
vertices. The red color means the subelements are found by rotating, whereas the blue ones need an additional mirroring.

\renewcommand{\arraystretch}{1.5}

\begin{figure}[bhtp!]
  \begin{minipage}[c]{0.45\linewidth}\centering
    \includegraphics[width=0.9\textwidth]{el1_tet}
  \end{minipage}
  \begin{minipage}[c]{0.45\linewidth}\centering
    \begin{tabular}{c | p{5.2cm} | c }
      %\hline
      \footnotesize{subel.} & \footnotesize{\texttt{coor(1)}}  & \footnotesize{conf.}  \\ 
      \hline%\hline
      1 & ($x_0$,                  $y_0$,                  $z_0$)                  )        & 1 \\% \hline
      2 & ($x_0 + \frac{h(1)}{2}$, $y_0$,                  $z_0$                   )        & 2 \\% \hline
      3 & ($x_0 + \frac{h(1)}{2}$, $y_0$,                  $z_0$                   )        & 3 \\% \hline
      4 & ($x_0 + \frac{h(1)}{2}$, $y_0$,                  $z_0$                   )        & 1 \\% \hline
      5 & ($x_0 + \frac{h(1)}{2}$, $y_0 + \frac{h(2)}{2}$, $z_0$                   )        & 4 \\% \hline
      6 & ($x_0 + \frac{h(1)}{2}$, $y_0 + \frac{h(2)}{2}$, $z_0$                   )        & 1 \\% \hline
      7 & ($x_0 + \frac{h(1)}{2}$, $y_0 + \frac{h(2)}{2}$, $z_0$                   )        & 5 \\% \hline
      8 & ($x_0 + \frac{h(1)}{2}$, $y_0 + \frac{h(2)}{2}$, $z_0 + \frac{h(3)}{2}$  )        & 1% \\% \hline
    \end{tabular}
  \end{minipage}
\caption{Subdivision of tetrahedrons with \texttt{conf = 1}}    
\label{fig:tet_el1}
\end{figure}



\begin{figure}[bhtp!]
  \begin{minipage}[c]{0.45\linewidth}\centering
    \includegraphics[width=1.0\textwidth]{el2_tet}
  \end{minipage}
  \begin{minipage}[c]{0.45\linewidth}\centering
    \begin{tabular}{c | p{5.2cm} | c }
      %\hline
      \footnotesize{subel.} & \footnotesize{\texttt{coor(1)}}  & \footnotesize{conf.}  \\ 
      \hline%\hline
      1 & ($x_0$,                  $y_0$,                   $z_0$)                  )        & 2 \\ %\hline
      2 & ($x_0$,                  $y_0 + \frac{h(2)}{2}$,  $z_0$                   )        & 2 \\ %\hline
      3 & ($x_0$,                  $y_0 + \frac{h(2)}{2}$,  $z_0$                   )        & 6 \\ %\hline
      4 & ($x_0$,                  $y_0 + \frac{h(2)}{2}$,  $z_0$                   )        & 4 \\ %\hline
      5 & ($x_0$,                  $y_0 + \frac{h(2)}{2}$,  $z_0 + \frac{h(3)}{2}$  )        & 2 \\ %\hline
      6 & ($x_0$,                  $y_0 + \frac{h(2)}{2}$,  $z_0 + \frac{h(3)}{2}$  )        & 3 \\ %\hline
      7 & ($x_0$,                  $y_0 + \frac{h(2)}{2}$,  $z_0 + \frac{h(3)}{2}$  )        & 1 \\ %\hline
      8 & ($x_0 + \frac{h(1)}{2}$, $y_0 + \frac{h(2)}{2}$,  $z_0 + \frac{h(3)}{2}$  )        & 2 %\\ %\hline
    \end{tabular}
  \end{minipage}
\caption{Subdivision of tetrahedrons with \texttt{conf = 2}}    
\label{fig:tet_el2}
\end{figure}




\begin{figure}[bhtp!]
  \begin{minipage}[c]{0.45\linewidth}\centering
    \includegraphics[width=1.0\textwidth]{el3_tet}
  \end{minipage}
  \begin{minipage}[c]{0.45\linewidth}\centering
    \begin{tabular}{c | p{5.2cm} | c }
      %\hline
      \footnotesize{subel.} & \footnotesize{\texttt{coor(1)}}  & \footnotesize{conf.}  \\ 
      \hline%\hline
      1 & ($x_0$,                   $y_0$,                   $z_0$)                  )        & 3 \\ %\hline
      2 & ($x_0$,                   $y_0 + \frac{h(2)}{2}$,  $z_0$                   )        & 5 \\ %\hline
      3 & ($x_0$,                   $y_0 + \frac{h(2)}{2}$,  $z_0$                   )        & 3 \\ %\hline
      4 & ($x_0$,                   $y_0 + \frac{h(2)}{2}$,  $z_0$                   )        & 1 \\ %\hline
      5 & ($x_0 + \frac{h(1)}{2}$,  $y_0 + \frac{h(2)}{2}$,  $z_0$                   )        & 6 \\ %\hline
      6 & ($x_0 + \frac{h(1)}{2}$,  $y_0 + \frac{h(2)}{2}$,  $z_0$                   )        & 2 \\ %\hline
      7 & ($x_0 + \frac{h(1)}{2}$,  $y_0 + \frac{h(2)}{2}$,  $z_0$                   )        & 3 \\ %\hline
      8 & ($x_0 + \frac{h(1)}{2}$,  $y_0 + \frac{h(2)}{2}$,  $z_0 + \frac{h(3)}{2}$  )        & 3 %\\ %\hline
    \end{tabular}
  \end{minipage}
\caption{Subdivision of tetrahedrons with \texttt{conf = 3}}    
\label{fig:tet_el3}
\end{figure}




\begin{figure}[bhtp!]
  \begin{minipage}[c]{0.45\linewidth}\centering
    \includegraphics[width=1.0\textwidth]{el4_tet}
  \end{minipage}
  \begin{minipage}[c]{0.45\linewidth}\centering
    \begin{tabular}{c | p{5.2cm} | c }
      %\hline
      \footnotesize{subel.} & \footnotesize{\texttt{coor(1)}}  & \footnotesize{conf.}  \\ 
      \hline%\hline
      1 & ($x_0$,                   $y_0$,                   $z_0$)                  )        & 4 \\% \hline
      2 & ($x_0$,                   $y_0$,                   $z_0 + \frac{h(3)}{2}$  )        & 1 \\% \hline
      3 & ($x_0$,                   $y_0$,                   $z_0 + \frac{h(3)}{2}$  )        & 5 \\% \hline
      4 & ($x_0$,                   $y_0$,                   $z_0 + \frac{h(3)}{2}$  )        & 4 \\ %\hline
      5 & ($x_0 + \frac{h(1)}{2}$,  $y_0$,                   $z_0 + \frac{h(3)}{2}$  )        & 6 \\% \hline
      6 & ($x_0 + \frac{h(1)}{2}$,  $y_0$,                   $z_0 + \frac{h(3)}{2}$  )        & 2 \\% \hline
      7 & ($x_0 + \frac{h(1)}{2}$,  $y_0$,                   $z_0 + \frac{h(3)}{2}$  )        & 4 \\% \hline
      8 & ($x_0 + \frac{h(1)}{2}$,  $y_0 + \frac{h(2)}{2}$,  $z_0 + \frac{h(3)}{2}$  )        & 4 %\\% \hline
    \end{tabular}
  \end{minipage}
\caption{Subdivision of tetrahedrons with \texttt{conf = 4}}    
\label{fig:tet_el4}
\end{figure}




\begin{figure}[bhtp!]
  \begin{minipage}[c]{0.45\linewidth}\centering
    \includegraphics[width=1.0\textwidth]{el5_tet}
  \end{minipage}
  \begin{minipage}[c]{0.45\linewidth}\centering
    \begin{tabular}{c | p{5.2cm} | c }
      %\hline
      \footnotesize{subel.} & \footnotesize{\texttt{coor(1)}}  & \footnotesize{conf.}  \\ 
      \hline%\hline
      1 & ($x_0$,                   $y_0$,                   $z_0$                   )        & 5 \\% \hline
      2 & ($x_0 + \frac{h(1)}{2}$,  $y_0$,                   $z_0$                   )        & 6 \\% \hline
      3 & ($x_0 + \frac{h(1)}{2}$,  $y_0$,                   $z_0$                   )        & 4 \\% \hline
      4 & ($x_0 + \frac{h(1)}{2}$,  $y_0$,                   $z_0$                   )        & 5 \\% \hline
      5 & ($x_0 + \frac{h(1)}{2}$,  $y_0$,                   $z_0 + \frac{h(3)}{2}$  )        & 3 \\% \hline
      6 & ($x_0 + \frac{h(1)}{2}$,  $y_0$,                   $z_0 + \frac{h(3)}{2}$  )        & 1 \\% \hline
      7 & ($x_0 + \frac{h(1)}{2}$,  $y_0$,                   $z_0 + \frac{h(3)}{2}$  )        & 5 \\% \hline
      8 & ($x_0 + \frac{h(1)}{2}$,  $y_0 + \frac{h(2)}{2}$,  $z_0 + \frac{h(3)}{2}$  )        & 5% \\% \hline
    \end{tabular}
  \end{minipage}
\caption{Subdivision of tetrahedrons with \texttt{conf = 5}}    
\label{fig:tet_el5}
\end{figure}

\begin{figure}[bhtp!]
  \begin{minipage}[c]{0.45\linewidth}\centering
    \includegraphics[width=1.0\textwidth]{el6_tet}
  \end{minipage}
  \begin{minipage}[c]{0.45\linewidth}\centering
    \begin{tabular}{c | p{5.2cm} | c  }
      %\hline
      \footnotesize{subel.} & \footnotesize{\texttt{coor(1)}}  & \footnotesize{conf.}  \\ 
      \hline%\hline
      1 & ($x_0$,                   $y_0$,                   $z_0$                   )        & 6 \\% \hline
      2 & ($x_0$,                   $y_0$,                   $z_0 + \frac{h(3)}{2}$  )        & 6 \\% \hline
      3 & ($x_0$,                   $y_0$,                   $z_0 + \frac{h(3)}{2}$  )        & 2 \\ %\hline
      4 & ($x_0$,                   $y_0$,                   $z_0 + \frac{h(3)}{2}$  )        & 3 \\ %\hline
      5 & ($x_0$,                   $y_0 + \frac{h(2)}{2}$,  $z_0 + \frac{h(3)}{2}$  )        & 6 \\% \hline
      6 & ($x_0$,                   $y_0 + \frac{h(2)}{2}$,  $z_0 + \frac{h(3)}{2}$  )        & 4 \\% \hline
      7 & ($x_0$,                   $y_0 + \frac{h(2)}{2}$,  $z_0 + \frac{h(3)}{2}$  )        & 5 \\% \hline
      8 & ($x_0 + \frac{h(1)}{2}$,  $y_0 + \frac{h(2)}{2}$,  $z_0 + \frac{h(3)}{2}$  )        & 6% \\% \hline
    \end{tabular}
  \end{minipage}
\caption{Subdivision of tetrahedrons with \texttt{conf = 6}}    
\label{fig:tet_el6}
\end{figure}


The only problem is that the reference region on which the eltree subdivision takes places (the Hill-simplex) is 
different from the reference region for isoparametric elements for tetrahedrons in \tfem. To solve this problem, a 
simple mapping is performed which is shown in Figure \ref{fig:mapping}. Note that the mapping is only defined for a 
tetrahedron of \texttt{conf = 2} (because this transformation is relatively easy), therefore, the root element must 
always have \texttt{conf = 2}

\begin{figure}[bhtp!]
\begin{center}
\includegraphics[width = 0.8\linewidth]{mapping}
\caption{Mapping from the eltree reference region to the isoparametric reference region for tetrahedrons.}
\label{fig:mapping}
\end{center}
\end{figure}

To fill the root element for tetrahedrons correctly, the first call to \texttt{fill\_node\_eltree} will always look like: \newline

\noindent
\texttt{coor(1,:) = [ 0.0\_dp, 0.0\_dp, 0.0\_dp ] } \newline
\texttt{coor(2,:) = [ 1.0\_dp, 1.0\_dp, 1.0\_dp ] } \newline
\texttt{conf = 2}\newline
\texttt{eltype = 1}\newline
\texttt{call fill\_node\_eltree ( eltree, coor, conf, eltype )}\newline
